% Include after sections/supplement/supplement_c_geometry in supplement_personnel.tex.
\clearpage
\subsection*{C16. Parameter, validation and implementation details}
\label{sec:consensus-implementation}



\paragraph{Quality-score development.}
Parameter development used controlled geometry and development cases. 28 configurations compared direction scales of approximately 6.12, 8, 10 and 12 degrees, boundary denominators of 32, 48 and 64, and combination exponents of 1.25, 1.5, 2 and 3. The retained \(8^\circ\) scale responds to isolated direction deviations; L2 moderates their accumulation with footprint loss, and the boundary denominator of 32 retains sensitivity to omissions and localization errors. Development comparisons of global-height scales 1.0 and 0.8 and local-height settings near 0.25 informed the retained height scales. The formula was then fixed for validation and personnel analyses.

\paragraph{Score-scale sensitivity.}
The sensitivity analysis uses 9 distinct parameter settings: the primary score and 8 variants multiplying one scale by $0.75$ or $1.25$. The $1.25\times$ direction variant is the $10^\circ$ setting. Boundary variants scale $\sqrt{32}$, giving squared denominators of 18 and 50; global-height scales are 0.6 and 1.0, local-height scales 0.1875 and 0.3125, and direction scales $6^\circ$ and $10^\circ$.

The 3--6-person comparison uses 206 images and the 20--30-person comparison uses 171. Each interval retains the same jointly scorable paths across settings within each image; P10 changes subtract endpoint quantiles. The 14-person matched replacement comparison uses 69 images in 20 buildings and retains the aggregation defined in C17. Every setting preserves early mean gains, late P10 gains, reduced late standard deviation and positive replacement excess (Table~\ref{tab:score-scale-sensitivity}). The existing parameter and result records accompany the manuscript in \path{supplement_sources/score_sensitivity/}.

\begin{table}[htbp]
\centering\small
\caption{Score-scale sensitivity on fixed output supports. All values are Q-point changes or Q-point excesses. Each row is a distinct parameter setting.}
\label{tab:score-scale-sensitivity}
\setlength{\tabcolsep}{4pt}
\begin{tabular}{@{}lrrrrr@{}}
\toprule
Setting & Mean 3--6 & Mean 20--30 & P10 20--30 & SD 20--30 & $G_S$ \\
\midrule
Primary & +4.738 & +0.502 & +2.844 & -2.005 & +3.441 \\
Global height $\times0.75$ & +4.743 & +0.503 & +2.844 & -2.003 & +3.463 \\
Global height $\times1.25$ & +4.735 & +0.502 & +2.846 & -2.006 & +3.430 \\
Local height $\times0.75$ & +4.832 & +0.515 & +2.865 & -2.009 & +3.455 \\
Local height $\times1.25$ & +4.692 & +0.497 & +2.835 & -2.003 & +3.432 \\
Direction $\times0.75$ & +5.223 & +0.679 & +3.251 & -2.205 & +3.460 \\
Direction $\times1.25$ ($10^\circ$) & +4.374 & +0.400 & +2.597 & -1.880 & +3.385 \\
Boundary scale $\times0.75$ & +5.011 & +0.510 & +2.999 & -2.104 & +3.989 \\
Boundary scale $\times1.25$ & +4.600 & +0.505 & +2.784 & -1.955 & +3.158 \\
\bottomrule
\end{tabular}
\end{table}


\paragraph{8-reviewer questionnaire protocol and defect relations.}
2 Manual annotations were randomly selected from each of 30 panoramas across 17 buildings: 9 easy, 9 moderate and 9 difficult ordinary panoramas plus 3 doorway panoramas. 8 annotators were randomly invited to take part in the questionnaire. They rated each annotation as good, medium or poor and selected visible defects, using only panorama overlays, shared-scale floor footprints and 3D views in a fixed randomized arrangement. Grades were encoded as 2, 1 and 0, respectively, before averaging.

Defect-selection shares use valid answers for each annotation: unknown and blank answers are excluded, while ``none'' is a valid report of no visible defect. The 480 defect answers comprise 460 valid reports and 20 unknown reports, with no blank answers or combinations of ``none'' with a specific defect. Across the 60 annotations, missing-structure shares correlate with missing effective turns (\(\rho=0.544\)) and floor omission (\(\rho=0.240\)); explicit floor-omission shares correlate with omitted floor fraction (\(\rho=0.586\)). Extra-structure shares correlate with extra turns (\(\rho=0.415\)) and floor extension (\(\rho=0.157\)). Boundary-position shares correlate with \(S_b\) (\(\rho=0.395\)) and \(D_{\mathrm{dir}}\) (\(\rho=0.358\)). Wall-height shares correlate with \(H_m\) (\(\rho=0.090\)), \(H_c\) (\(\rho=0.612\)) and the combined height loss (\(\rho=0.609\)), distinguishing local height discrepancies from an overall height shift.

Of 480 submitted grades, 478 are valid (67 good, 163 medium and 248 poor); 2 are unknown. Within-reviewer direction compares the grades of the two annotations on the same panorama; panorama-level direction averages the common valid reviewers first. Reviewer agreement counts equal grades among valid reviewer pairs for each annotation. Correlation intervals use 10,000 bootstrap resamples of the 17 buildings, retaining all their images, paired annotations and reviewers; sensitivity omits each reviewer in turn. Supplement D5 reports these additional grade-based results.

The questionnaire responses, per-annotation geometric measurements and recomputation scripts are provided in \path{supplement_sources/questionnaire/}.

\paragraph{Additional geometric review.}
A separate geometric review examined 25 same-image pairs: 23 Manual and 2 Semi. 20 pairs were selected for specified geometric conditions and score relationships; 5 mechanism examples were already familiar to the reviewer. Across all 25 pairs, 10 received directional judgments, 3 were judged close, and 12 were difficult to compare or problematic on both sides. Among the 20 condition-selected pairs, 6 judgments were directional. Q with direction scales \(8^\circ\) and \(10^\circ\) agreed with all 5 directional Manual judgments, while prism IoU agreed with 3. In the directional Semi pair, the reviewer preferred the lower-scoring response.

\paragraph{Working skip penalty.}
The matching uses a fixed penalty of \(g=7.5^\circ\) for each unmatched observation. Selected through tuning, this value accommodates positional differences at the same wall turn while keeping nearby distinct turns separate.

\paragraph{Candidate identity groups.}
Pairwise cyclic alignment uses the average spherical distance of the upper
and lower endpoints and a skip penalty of \(g=7.5^\circ\) for each unmatched
observation. Complete-linkage agglomerative clustering encodes these
alignments with dissimilarity zero on the diagonal, \(c(a,b)/180^\circ\) for
observations selected as a pair in the corresponding worker-pair alignment,
and 2 otherwise. The hierarchy is cut at 1. Each resulting candidate group
therefore contains at most one observation per worker, with matching support
for every pair of observations. Pairwise links through an intermediate
observation alone are insufficient to join a group.

\paragraph{Periodic positioning and ties.}
The observed-pair anchor minimizes the sum of the larger upper- or
lower-endpoint angular distances to the other group members. Anchor ties
follow the fixed worker, response and pair-index order. With panorama
width \(W=1024\), anchor coordinate \(x_0\), and wrapped offsets
\(d_j=((x_j-x_0+W/2)\bmod W)-W/2\), a position is returned only if
\(\max_j d_j-\min_j d_j<W/2\) and no \(|d_j|=W/2\). The boundary checks
use a \(10^{-9}\)-pixel tolerance. A failed check leaves the position
unresolved. Horizontal and vertical medians use the same group's
observations with equal worker weights; their resulting coordinates may
lie between original observations.

\paragraph{Collective geometry.}
The retained upper--lower pairs define the collective wall layout, which undergoes the geometric validity checks in C15. Team records preserve membership, source observations, correspondence groups, support counts and aggregated coordinates.

\paragraph{Floor-region comparison and local structure.}
BEV voting partitions the team's overlapping floor polygons into regions and retains regions covered by at least half the people. Its output is a two-dimensional floor region. Paired-corner consensus retains the upper and lower coordinates of supported wall turns and connects them into a wall layout.

Across 138 teams from 48 pools, BEV voting yields 138 valid floor footprints; paired-corner consensus yields 131 footprints and 130 complete 3D layouts. Of the 70 standard Manual teams, 67 yield complete evaluable layouts. Across these 67 layouts, consensus retains 354 strictly majority-supported and 28 half-supported clusters, while excluding 116 minority clusters.

Figure~\ref{fig:consensus-unb45} shows the full panorama and the corresponding floor outputs for a 10-person team. The same 7 people mark E and F. Both the paired-corner layout and the BEV region retain the resulting local change in floor coverage; the paired-corner output also records the supporters and upper--lower positions of those turns.

\begin{figure}[htbp]
\centering
\includegraphics[width=\linewidth]{supplement_figures/consensus_unb45.pdf}
\caption{Local wall turns enter a collective layout through their supporters. In panorama uNb9QFRL6hY-45, E and F each receive support from the same 7 of 10 annotators. The paired-corner layout retains 6 turns, while BEV voting retains the corresponding floor-region change. The panorama and close view locate E and F; floor views share coordinates and scale, with the reference shown by dashed lines.}
\label{fig:consensus-unb45}
\end{figure}
